\begin{claim}{D.2.34}{Proved}{The remaining Toda value vanishes}
Keep the spaces and the normalization of \textup{D.2.33}. Then
\[
 \Delta=\langle\lambda,\eta_8,2\rangle
             =\{0\}\subseteq\pi_{10}W.
 \tag{D.2.34a}
\]
Consequently every ordinary pointed map
$h:\Sigma A\to Y$ has $\epsilon=0$ in the normalized
once-suspended functional of \textup{D.2.28}.
The simultaneous first-cone vanishing of \textup{D.2.27}
therefore holds for every mixture of suspension exponents
$0,1,2,3$ and at least $5$. Exponent $4$ and propagation
through later cones are not settled by this theorem.
\end{claim}
\begin{proof}
Write $\rho_T:S^6\to T$ for the actual sphere in
\textup{D.2.17}, and $\rho_Y:S^8\to Y$ for its double
adjoint. With the homology labels used there,
\[
 h_2(\rho_T)=b^2+\{a,b\}+f+ac+a^3,
       \qquad h_2(\rho_Y)=\sigma^2f.
 \tag{D.2.34b}
\]
We first identify an actual representative of the shifted
generator $\lambda_Y\in\pi_9Y\cong\ZZ/2$.
Put
\[
 h_0=\eta_2\wedge\id_B:\Sigma Y\longrightarrow Y,
       \qquad \varphi_0:Y\longrightarrow\Omega Y
 \tag{D.2.34c}
\]
for this map and its once adjoint. The quadratic calculation
of \textup{D.2.13} applies in degree eight and gives
\[
 \varphi_{0*}(\sigma^2f)
       =(\sigma a)(\sigma c)+(\sigma c)(\sigma a).
 \tag{D.2.34d}
\]
Here is the complete degree check. In the cone
$\mathbb CP^2\wedge B$, the product of the two old
classes $tu$ and $tu^2$ is $t^2u^3$. Thus the reduced
diagonal of the corresponding new class has precisely
the two ordered terms giving \textup{(D.2.34d)}.
Equivalently, $\overline\Delta f=a\otimes c+c\otimes a$.
The linear component vanishes because $h_0$ induces zero
on reduced homology. In the James tensor algebra
$H_*(\Omega\Sigma^2B;\FF_2)$, the bottom label
$\sigma a$ has degree three, so tensor length three
starts in degree nine and cannot contribute here.
The relative Adams--Hilton formula for this quadratic
coefficient is the one already audited in
\cite[\S3.1, p.~300; proof of Proposition~7.4,
p.~309]{FGHHalperinLemaire1995}; the tensor-algebra
identification is \cite[\S3; Theorem~5.6 and
Corollary~5.7]{I_James1955}.

The right side of \textup{(D.2.34d)} is the nonzero
homology suspension of $R=\{a,c\}$, which detects the
generator in \textup{D.2.19}. Applying this calculation
to \textup{(D.2.34b)} proves the equality of actual
homotopy classes
\[
                 \lambda_Y=h_0\circ\Sigma\rho_Y.
 \tag{D.2.34e}
\]
No formula on higher nonprimitive labels is used.

The relation $\rho_T\eta_6=0$ is part of
\textup{(D.2.19a)}. Choose an actual extension
\[
 \overline\rho:C\eta_6=S^6\cup_{\eta_6}e^8
                      \longrightarrow T.
\]
The ordinary cofiber exact sequence gives
\[
 \pi_8(C\eta_6)=\ZZ_{(2)}\{\Gamma\},
       \qquad h_{\ZZ}(\Gamma)=2[e^8].
 \tag{D.2.34f}
\]
Indeed the preceding multiplication by $\eta$ from
$\pi_1^S$ to $\pi_2^S$ is an isomorphism, and the
image in the integral top-cell group consists exactly
of the even integers. These are ordinary groups in
the stable range of the five-connected space $C\eta_6$;
the sphere identities are those in
\cite[\S2, (2.1), p.~149]{I_Oka1986}.
The class $\Gamma$ is a coextension of degree two
across $\eta_6$, and its mod-two Hurewicz image is zero.
Consequently
\[
 h_2(\overline\rho\Gamma)=0.
\]
The decisive input is the injectivity of the actual
ordinary Hurewicz map
\[
               \pi_8T\longrightarrow H_8(T;\FF_2)
\]
proved in \textup{D.2.20}. It gives
$\overline\rho\Gamma=0$, an actual zero representative
of $\langle\rho_T,\eta_6,2\rangle$.

Double suspension followed by the double-loop counit
transports this extension, this coextension, and this
nullhomotopy to a zero representative of
$\langle\rho_Y,\eta_8,2\rangle$ in $\pi_{10}Y$.
Suspending once more and postcomposing with $h_0$,
the extension/coextension definition of the Toda bracket
and \textup{(D.2.34e)} give
\[
            0\in\langle\lambda_Y,\eta_9,2\rangle
                    \subseteq\pi_{11}Y.
 \tag{D.2.34g}
\]
This uses a specified zero representative, so no
surjectivity assertion about suspension of bracket sets
is required. All the cone domains here have bottom cell
above degree four; their maps and homotopies lift through
$Z=\fib(Y\to K(\FF_2,4))$ with the needed uniqueness.
Adjunction therefore gives a zero representative of the
bracket in \textup{(D.2.34a)} in $W=\Omega Z$.
Its full indeterminacy is zero by
\textup{(D.2.33o)}, using \textup{D.2.31}; hence that
bracket is the singleton $\{0\}$.

Finally \textup{D.2.33} shows that any ordinary map with
$\epsilon=1$ would force $\Delta=p$, where $p\ne0$.
Equation \textup{(D.2.34a)} excludes this possibility
for every such map, proving the stated uniform conclusion.
The normalization is the same single target class used in
\textup{D.2.27--D.2.28}. The common-factor criterion of
\textup{D.2.26} therefore combines this vanishing with all
the exponents already treated in \textup{D.2.27}, including
arbitrary simultaneous mixtures of them.
\end{proof}
